%%
%% Introduction
%%
\section{Introduction}
\label{sec:introduction}
Matrix operations form the computational backbone of modern machine learning. From the weighted sums in neural network layers to the attention mechanisms in transformers, nearly every inference and training step reduces to matrix-vector or matrix-matrix multiplication~\cite{alphatensor}. As models scale to billions of parameters, these operations become the dominant computational bottleneck, motivating extensive research into hardware accelerators and algorithmic improvements.

Quantum computing offers a fundamentally different approach to these computations. Algorithms such as HHL for linear systems~\cite{hhl} and quantum machine learning models~\cite{qml-nature} leverage quantum parallelism to achieve potential speedups for linear algebra workloads. As quantum hardware matures, the ability to express and compile such computations efficiently becomes increasingly important.

For researchers and developers in machine learning and scientific computing, however, this creates a significant adoption barrier. Research teams and organizations typically maintain large codebases in C or C++ containing optimized numerical routines developed over years. Rewriting these algorithms in low-level quantum assembly is impractical and error-prone. A high-level compiler that translates standard C into quantum circuits addresses this gap: it reduces integration costs by allowing developers to selectively target compute-intensive subroutines---such as matrix multiplications in training loops or inference kernels---for quantum acceleration, without abandoning existing code or familiar development workflows.

However, current high-level quantum compilers primarily support scalar integer arithmetic~\cite{quantumc,qhls}, leaving a significant gap in expressiveness for developers who wish to leverage quantum computing for linear algebra workloads. While these compilers successfully demonstrate the feasibility of translating classical code into quantum circuits, they cannot directly handle the vector and matrix operations that dominate modern computational workloads. This limitation forces developers to either manually decompose their algorithms into scalar operations or abandon high-level compilation altogether.

\textbf{Related Work.}
Several approaches have been proposed to raise the abstraction level of quantum programming. Quipper~\cite{quipper} provides a scalable functional framework embedded in Haskell, enabling the implementation of complex algorithms such as those in the IARPA QCS program. Silq~\cite{silq} introduces a high-level quantum language with automatic uncomputation through a novel type system, significantly simplifying the management of temporary qubits. ScaffCC~\cite{scaffcc} compiles the Scaffold language to quantum circuits using LLVM infrastructure, providing mature compiler optimizations for quantum programs. While these approaches successfully raise the abstraction level, they all require developers to adopt new programming languages or paradigms.

A complementary line of work focuses on compiling standard imperative languages to quantum circuits. QuantumC~\cite{quantumc} introduced an MLIR-based pipeline that translates C programs with scalar integer operations into OpenQASM circuits. QHLS~\cite{qhls} explores a similar C-to-quantum synthesis approach but relies on ad-hoc transformations instead of a modular IR infrastructure. Both systems demonstrate the feasibility of high-level quantum compilation from C, but neither supports array, vector, or matrix operations---the computational primitives that dominate machine learning and scientific computing workloads.

At the intermediate representation level, QIRO~\cite{qiro} proposes an SSA-based quantum program representation designed for optimization, demonstrating the benefits of leveraging classical compiler techniques in the quantum domain. Our work similarly builds on SSA semantics through MLIR, but addresses the front-end challenge of recognizing and compiling structured linear algebra computations from C source code.

\textbf{This Work.} This work extends the QuantumC compilation framework to address a fundamental limitation of existing high-level quantum compilers: the lack of native support for vector and matrix computations. Rather than lowering array-based programs through scalar loop unrolling, we adopt a semantics-preserving approach that explicitly identifies linear algebra computations at the source level. By recognizing structured access patterns and loop constructs in C programs, the compiler can lift vector and matrix operations into dedicated intermediate representations, retaining high-level information throughout the compilation flow.

This design enables the compilation of hybrid programs that combine scalar control logic with vector and matrix arithmetic, while preserving compatibility with the original QuantumC pipeline. In addition, the extended framework supports compile-time initialization of array data, allowing constant vectors and matrices to be directly encoded into the generated quantum circuits. Beyond expressiveness, this extension provides a concrete basis for analyzing different arithmetic realization strategies, which we exploit to compare alternative quantum adder implementations within a unified compilation setting.

\textbf{Contributions.} This work extends the QuantumC compiler with native support for vector and matrix computations. We introduce a pattern-based recognition mechanism that identifies three linear algebra patterns---vector addition, dot product, and matrix multiplication---directly from C source code, lifting them into dedicated intermediate representations. Unlike naive loop unrolling, this approach preserves the inherent parallelism of vector operations: experimental results show that circuit depth remains constant regardless of vector size (depth 23 for both $N=2$ and $N=4$), whereas sequential unrolling would produce depth linear in $N$. The extended pipeline also supports hybrid programs that combine scalar control logic with vector and matrix arithmetic. Finally, we leverage this unified infrastructure to provide a comparative analysis of QFT-based and ripple-carry arithmetic backends, characterizing their trade-offs in terms of circuit depth, gate composition, and qubit requirements across different workload types.

\textbf{Structure of the paper.} The remainder of this paper is organized as follows. Section~\ref{sec:background} provides background on QuantumC and quantum arithmetic implementations. Section~\ref{sec:pipeline} describes our extended compilation pipeline, including pattern recognition and the VecMat/QAR intermediate representations. Section~\ref{sec:results} presents experimental results on vector and matrix benchmarks and the arithmetic backend comparison. Section~\ref{sec:conclusion} concludes and discusses future work.
